% PAGES: 224-225

% CONTINUES: closes the \begin{proof} of Theorem 7.3 left open at the end of
% sec07_c2_3.tex (PDF page 223, folio 207), whose last item was eq. (7.80).
If $\bfX = (X_i)_{i=1:n}$ denotes the column vector of the random variables $X_i$,
$i = 1:n$, then formula (7.24) from Lemma 7.3 can be written as follows:
\begin{equation}
\cov\!\left((C^{(1)})^t\bfX,(C^{(2)})^t\bfX\right) \;=\; (C^{(1)})^t\Sigma_{\bfX}C^{(2)} \;=\; (C^{(2)})^t\Sigma_{\bfX}C^{(1)}.
\end{equation}
Then, from (7.80) and using (7.81) it follows that
\begin{align}
\Sigma_{\bfY}(j,k) &= \cov(Y_j,Y_k) \;=\; \cov(r_j\bfX,r_k\bfX) \notag \\
&= r_j\Sigma_{\bfX}r_k^t, \quad \forall\, 1\le j,k\le m.
\end{align}
From (7.79) and (7.82), we conclude that $\Sigma_{\bfY} = M\Sigma_{\bfX}M^t$.
\end{proof}

\begin{examplex}
As a simple application of the Linear Transformation Property, let $X_1$,
$X_2$, \ldots, $X_n$ be nonconstant random variables, and let $Y_1$, $Y_2$,
\ldots, $Y_n$ be the random variables given by $Y_i = d_iX_i$, where $d_i \neq 0$,
are constants, for $i = 1:n$. Let $\bfY = (Y_i)_{i=1:n}$ and $\bfX =
(X_i)_{i=1:n}$. Then, $\bfY = D\bfX$, and, from the Linear Transformation Property
(7.76), it follows that $\Sigma_{\bfY} = D\Sigma_{\bfX}D^t$. Since $D$ is a
diagonal matrix, $D^t = D$, we conclude that
\begin{equation}
\Sigma_{\bfY} \;=\; D\Sigma_{\bfX}D.
\end{equation}
\end{examplex}

\begin{theorem}
Let $X_1$, $X_2$, \ldots, $X_n$ be $n$ random variables, and let $\bfX =
(X_i)_{i=1:n}$. Denote by $\mu_{\bfX}$ the mean vector of $\bfX$ and by
$\Sigma_{\bfX}$ the covariance matrix of $\bfX$. Let $Y_1$, $Y_2$, \ldots, $Y_m$ be
random variables given by
\[
\bfY \;=\; b + M\bfX,
\]
where $\bfY = (Y_i)_{i=1:m}$, $M$ is an $m \times n$ matrix, and $b$ is an $m
\times 1$ column vector. Denote by $\mu_{\bfY}$ the mean vector of $\bfY$ and by
$\Sigma_{\bfY}$ the covariance matrix of $\bfY$.

Then,
\begin{equation}
\mu_{\bfY} \;=\; b + M\mu_{\bfX} \qquad \text{and} \qquad \Sigma_{\bfY} \;=\; M\Sigma_{\bfX}M^t.
\end{equation}
\end{theorem}

\begin{proof}
From the linearity of expectation, it follows that
\begin{equation}
\E{M\bfX} \;=\; M\E{\bfX} \;=\; M\mu_{\bfX};
\end{equation}
see (7.128). Then, from (7.85), we find that
\[
\mu_{\bfY} \;=\; \E{\bfY} \;=\; \E{b+M\bfX} \;=\; b + \E{M\bfX} \;=\; b + M\mu_{\bfX}.
\]
Moreover, since $b$ is a constant vector, the covariance matrix of $\bfY = b +
M\bfX$ is the same as the covariance matrix of $M\bfX$, and therefore
$\Sigma_{\bfY} = \Sigma_{M\bfX}$. Note that $\Sigma_{M\bfX} = M\Sigma_{\bfX}M^t$;
see (7.77). Thus, we conclude that
\[
\Sigma_{\bfY} \;=\; \Sigma_{M\bfX} \;=\; M\Sigma_{\bfX}M^t.
\]
\end{proof}

The Linear Transformation Property can be used for generating normal random
variables with a given correlation matrix, which can be subsequently used for
Monte Carlo simulations, see section 7.5, and for establishing that any symmetric
positive semidefinite matrix is the covariance matrix (or the correlation matrix,
if the main diagonal entries are equal to $1$) of some random variables; see
section 7.4.

\section{Necessary and sufficient conditions for covariance and correlation matrices}

In Lemma 7.4, we proved that any covariance matrix is symmetric positive
semidefinite. We now show, using the Linear Transformation Property, that this is
a necessary and sufficient condition by finding normal random variables with
covariance matrix equal to any given symmetric positive semidefinite matrix, and
establish a similar result for correlation matrices.

\begin{theorem}
(i) An $n \times n$ square matrix is the covariance matrix of $n$ random variables
if and only if the matrix is symmetric positive semidefinite.

(ii) An $n \times n$ square matrix is the correlation matrix of $n$ random
variables if and only if the matrix is symmetric positive semidefinite and has
all the entries on the main diagonal equal to $1$.
\end{theorem}

\begin{proof}
(i) Recall from Lemma 7.4 that any covariance matrix is symmetric positive
semidefinite.

To prove that any symmetric positive semidefinite matrix is a covariance matrix,
let $A$ be an $n \times n$ symmetric positive semidefinite matrix. We will find
random variables $X_1$, $X_2$, \ldots, $X_n$ with covariance matrix
$\Sigma_{\bfX} = A$.

Since the matrix $A$ is symmetric, it follows from Theorem 5.4 that there exists
an orthogonal matrix $Q$ and a diagonal matrix $\Lambda$ such that
\begin{equation}
A \;=\; Q\Lambda Q^t.
\end{equation}
Recall that $\Lambda = \diag(\lambda_i)_{i=1:n}$, where $\lambda_i$, $i = 1:n$,
are the eigenvalues of $A$. Note that $\lambda_i \ge 0$, for all $i = 1:n$, since
the eigenvalues of a symmetric positive semidefinite matrix are nonnegative; see
Theorem 5.6. Let $\Lambda^{1/2}$ be the diagonal matrix given by
\begin{equation}
\Lambda^{1/2} \;=\; \diag\!\left(\sqrt{\lambda_i}\right)_{i=1:n}.
\end{equation}
Using (1.93), we find that
\begin{equation}
\Lambda^{1/2}\Lambda^{1/2} \;=\; \diag(\lambda_i)_{i=1:n} \;=\; \Lambda.
\end{equation}
Let
\begin{equation}
M \;=\; Q\Lambda^{1/2}.
\end{equation}
Then,
\begin{equation}
M^t \;=\; \left(\Lambda^{1/2}\right)^tQ^t \;=\; \Lambda^{1/2}Q^t,
\end{equation}
since $\Lambda^{1/2}$ is a diagonal matrix and therefore symmetric, i.e.,
$\left(\Lambda^{1/2}\right)^t = \Lambda^{1/2}$. Then, from (7.89), (7.90), (7.88),
and (7.86), we obtain that
\begin{equation}
MM^t \;=\; Q\Lambda^{1/2}\Lambda^{1/2}Q^t \;=\; Q\Lambda Q^t \;=\; A.
\end{equation}
Let $Z_1$, $Z_2$, \ldots, $Z_n$ be independent standard normal variables, and let
$X_1$, $X_2$, \ldots, $X_n$ be random variables\footnote{Note that $X_1$, $X_2$,
\ldots, $X_n$ are, in fact, normal random variables, since they are linear
combinations of independent normal variables.} given by $\bfX = M\bfZ$, where $M$
is the matrix given by (7.89), $\bfX = (X_i)_{i=1:n}$ and $\bfZ = (Z_i)_{i=1:n}$.
Let $\Sigma_{\bfX}$ be the covariance matrix of $X_1$,
% CONTINUES: \begin{proof} of Theorem 7.5 part (i) (opened above) is left OPEN
% here at the end of PDF page 225 (folio 209), mid-sentence right after "$X_1$,".
% The next file (sec07_c2_5.tex, PAGES 226) picks up with "$X_2, \ldots, X_n$."
% and closes this proof. Do not open a new \begin{proof} there.
